\section{Text watermarking: generation, removal and preservation}
\label{sec:text_watermarks}
\paragraph{Question and access.} A watermark differs from a content detector because its statistical signal is deliberately inserted using an owner-controlled mechanism. This part compares a released sampler watermark with a learned in-model adaptation. The reader can start here without reading the detector benchmarks or reinforcement-learning experiments.

\paragraph{Two distinct watermark locations.} TextSeal 2 marks the next-token sampling process while leaving the underlying Qwen weights unchanged. The DistSeal-inspired text adaptation first trains a hidden-state embedder and text extractor, then distils teacher distributions into Qwen LoRA weights and merges the student. The latter is a local experimental adaptation of the latent teacher/student idea, not an official implementation of a published DistSeal text system. Its teacher trains for 7,500 steps, and the rank-128 student for 8,000 steps. Teacher validation bit recovery reaches 70.36\%; that quantity is not the fraction of fresh student responses detected.

\paragraph{Data and generation.} Both arms use Qwen3.5-4B instruct with its native chat format and thinking disabled. HC3 full questions, rather than human-answer prefixes, form prompts. Normalized question hashes separate fitting, removal, validation and test pools; duplicate questions and answer prefixes are excluded. There are 256 test questions per arm. The recorded sampling configuration uses temperature one, nucleus sampling at cumulative probability 0.95, and a 200-token cap. Nucleus sampling keeps the most likely tokens needed to reach that mass, then renormalizes their probabilities.

\paragraph{The removal methods.} Ordinary rewriting and BIRA use an independent Llama-3.1-8B-Instruct rewriter with text-only access. Direct top-10 resampling and confidence smoothing instead generate fresh answers from simulated probability access. For TextSeal 2 this access reveals underlying pre-watermark logits. Sampling directly from those logits bypasses the external keyed sampler; it does not demonstrate the same bypass when a provider reveals only the keyed distribution. For the learned student, probabilities come from its merged marked weights. Smoothing combines target and independent Qwen3.5-0.8B probabilities using the adapted confidence rule.

\paragraph{Detection is not calibrated identically.} TextSeal 2 uses the released v2 deduplicated Gamma detector at nominal $p\leq .001$ and, secondarily, $p\leq .01$, where $p\in[0,1]$ is the null-tail $p$-value. The student uses extractor thresholds fitted on disjoint validation negatives. Its primary score threshold is $>-0.6263284683$. Comparing the two primary rates therefore compares two different operating-point constructions. A target-FPR label must not be presented as a precisely verified false-positive guarantee from only 256 null outputs.

\begin{table}[t]\centering\small
\begin{tabular}{lllrr}\toprule
Scheme & Condition & Access & Detected & Rate\\\midrule
TextSeal 2 & Original marked & -- & 253/256 & 98.8\% \\
TextSeal 2 & Clean control & -- & 1/256 & 0.4\% \\
TextSeal 2 & Ordinary rewrite & O & 15/256 & 5.9\% \\
TextSeal 2 & BIRA & O & 0/256 & 0.0\% \\
TextSeal 2 & Top-10 resampling & P & 1/256 & 0.4\% \\
TextSeal 2 & Smoothing & P & 0/256 & 0.0\% \\
Text student & Original marked & -- & 118/256 & 46.1\% \\
Text student & Clean control & -- & 2/256 & 0.8\% \\
Text student & Ordinary rewrite & O & 9/256 & 3.5\% \\
Text student & BIRA & O & 1/256 & 0.4\% \\
Text student & Top-10 resampling & P & 102/256 & 39.8\% \\
Text student & Smoothing & P & 18/256 & 7.0\% \\
\bottomrule\end{tabular}
\caption{Text-watermark experiments with Qwen3.5-4B instruct on 256 HC3 questions per condition. TextSeal 2 is the native keyed sampler; Text student is the local distilled latent-text adaptation. Detected is the number still flagged by each frozen owner test, and Rate is that count divided by 256. TextSeal 2 uses native $p\leq0.001$; the local student uses extractor score $>-0.6263284683$. Ordinary rewriting and BIRA use Llama-3.1-8B-Instruct on a saved marked answer. Top-10 resampling and smoothing generate fresh answers under simulated probability access; TextSeal 2 exposes pre-watermark probabilities, while the student exposes its marked model probabilities. Access: O means saved outputs only; P means a simulated probability-only target API without target weight access; -- marks a control. Initial detection differs, so post-attack detection alone does not count removals from initially detected sources.}\end{table}


\paragraph{The principal result.} TextSeal 2 begins with 253/256 detections. Ordinary rewriting leaves 15 detections and BIRA leaves none. The learned student begins with only 118/256 detections, so its post-attack rates alone can overstate removal: some outputs were already negative. The portable records additionally track paired transitions among initially positive examples. Direct resampling leaves 102/256 student detections, compared with 1/256 for TextSeal 2 under the stronger underlying-logit oracle.

\paragraph{Quality and semantic preservation.} Ordinary rewriting has mean Jina cosine 0.9339 and bidirectional NLI acceptance 79.7\% on TextSeal 2. BIRA has cosine 0.9108 and bidirectional NLI acceptance 47.3\%. A high cosine is thus compatible with an NLI model finding that the rewrite adds, drops or changes information. Numeric-set overlap is another diagnostic but does not check which number belongs to which entity or fact. Independent Mistral perplexity measures model likelihood, not truth.

A new post-hoc audit joins owner decisions, quality records and Jina scores by exact prompt ID. Its illustrative strict proxy requires bidirectional entailment, no NLI truncation, numeric-set F1 equal to one, valid perplexity, no wrapper leakage or long letter runs, and cosine at least .90. Missing or duplicate IDs cause an error. The audit does not tune the original attacks or thresholds.

\begin{table}[ht]\centering\small
\begin{tabular}{llrrrr}\toprule
Scheme & Rewrite & \shortstack{Initially\\detected} & Removed & \shortstack{Text-rule\\pass} & Joint\\\midrule
TextSeal 2 & Ordinary & 253 & 238 & 154 & 141 \\
TextSeal 2 & BIRA & 253 & 253 & 62 & 62 \\
Text student & Ordinary & 118 & 110 & 152 & 60 \\
Text student & BIRA & 118 & 117 & 62 & 28 \\
\bottomrule\end{tabular}
\caption{Post-hoc paired text-watermark audit on Qwen3.5-4B instruct, 256 source answers per row. TextSeal 2 uses native $p\leq0.001$; the local distilled student uses extractor score $>-0.6263284683$. Initially detected counts marked sources passing that test. Removed counts their attacked answers that become undetected. Text-rule pass counts all rewrites satisfying bidirectional NLI entailment without truncation, exact numeric-set agreement, valid positive independent Mistral perplexity, no wrapper leakage or long letter runs, and Jina cosine $\geq0.90$. Joint counts the intersection of removal and text-rule pass. Ordinary and BIRA rewriting both use independent Llama-3.1-8B-Instruct with output-only target access. These counts use automatic proxies and are not human annotations.}\end{table}


The strict proxy leaves 62 BIRA successes among 253 initially detected TextSeal 2 responses, versus 141 ordinary-rewrite successes. This changes the comparison substantially from raw detector suppression. These results are exploratory because the retained test outputs have already been inspected. The proxy is not a human-certified preservation criterion, and its cosine threshold is not a universal semantic constant. Confirmation experiments must freeze their criteria before producing new test outputs.

\paragraph{Length limits are part of the result.} In the retained TextSeal 2 arm, 254/256 original responses reach the 200-token budget. Some source answers are therefore incomplete, and a rewrite may finish an unfinished thought rather than preserve exactly the same claims. Some NLI pairs also exceed the evaluator's 512-token context. The full records retain these cases. A future quality-controlled experiment should distinguish complete responses, truncation, factual coverage and raw detectability rather than filtering inconvenient outputs after seeing their detector scores.

